\section{Custody and state compulsion}
\label{sec:14.3}
A government needn't alter a model to defeat a constraint; it can compel whoever holds it. The Apple case tested that in about as protected a jurisdiction as exists: a magistrate ordered Apple under the All Writs Act to write and sign software defeating passcode protections on a phone \autocite{pym2016applefbi}. Apple refused and moved to vacate, arguing the authority wouldn't stop at one phone. The question was never decided — the FBI found a third party and withdrew. So whether a democracy with an independent judiciary can force an unwilling company to compromise a system it built is legally untested. A government doesn't need that argument to lean on a vendor. The 2026 vendor dispute is the case: a vendor that refused to permit fully autonomous targeting and mass domestic surveillance was threatened with contract termination and designated a supply-chain risk. China's 2023 generative-AI rules exercise the same leverage as a standing condition of operating \autocite[arts.~4, 14, 17, 19, 21]{cac2023genai}. One government compelled case by case, against a vendor that refused, and the courts divided by statute, one setting aside the designation as retaliation and the other upholding a separate exclusion as turning on the refused contract term; the other compels in advance, as a term of doing business.

The American instrument that compels in advance is older than either case. Under Title I of the Defense Production Act the government places a rated order, which a supplier that normally makes the item must accept and fill ahead of other work; the grounds for rejecting one are narrow, objection to the use is not among them, and refusal carries a statutory penalty with no court in between \autocite{dpa1950,dpas2026}. A rated order says nothing about use. It compels delivery, and whether what the government does with the delivery is lawful stays the law of armed conflict's question, which leaves the vendor in the responsibility gap with the supply compelled. Two things follow. A vendor's refusal of a rated order is a refusal of a federal order at a price fixed in advance, which is the form that counts, and it is the form the 2026 dispute would have taken had the government carried out its threat to invoke the older statute \autocite{Amodei2026DoWStatement}. And a rated order gets the vendor's product. If the product holds the floor, the order gets a floored model, and compelling the vendor to remove the floor is compelling it to write code it doesn't want to write, which is the Apple question again, one step later and as untested.

Custody can be distributed without disappearing. Published weights can't be recalled by an order aimed at the developer; maintainers across jurisdictions raise the cost. None of that removes the leverage — a developer with employees, offices and legal existence inside a border answers to that government's courts regardless of how the model is built, and no architectural choice changes who can be held in contempt. Nobody has tested that leverage against a model built specifically to resist the government making the request.
